% Isolated after CP34 canonical cutoff; no canonical/shared input mutation.
\subsection{Residual modules on compatible cyclic source slices}

All geometric images and local systems below are taken over an
algebraically closed field of characteristic two. For every odd
cubic-principal pullback, its canonical proper submodule means
the unique maximal proper submodule supplied by the proved
full-hyperplane inertia and full-radical structure theorems.
It is defined without assuming that the image is full.

\begin{lemma}[A source projection lands in the source residual module]
\label{lem:DE-source-residual-projection}
\status{Candidate; complete proof supplied for independent audit.}
Let \(\eta:B\to C\) and \(\pi:A\to B\) be etale isogenies
of positive odd degrees \(a,n\), with \((C,\Theta_C)\)
nonordinary and principally polarized. Put
\[
 M=\eta^*\mathcal O_C(3\Theta_C),\quad L=\pi^*M,
 \quad h=9a.
 \tag{DE-CR1}
\]
Write \(V_0\) for the rank-\((h-1)\) source Lang space,
\(V\) for the rank-\((hn-1)\) target space, and \(K_0,K\)
for their canonical proper submodules. On a nonempty smooth,
sharp source-slice open, the actual outer deck projector \(e\)
and its splitting satisfy
\[
        V=V_0\oplus V_{\mathrm{nt}},\qquad
        e(V)=V_0,\qquad e(K)\subseteq K_0.
 \tag{DE-CR2}
\]
No full-image assumption is made on the source.
\end{lemma}
\begin{proof}
The first lemma of the arbitrary odd-source theorem supplies
the same nonempty extra-block determinant open and the same
actual composite common-wall points used in its conditional
quotient-bound proof. Thus the actual relative Lang splitting
and the regular-wall comparison in that proof hold unchanged,
before the step where a full source image was imposed.
Explicitly, the odd outer deck average is an idempotent on the
\(\mathbf F_2\)-local system. Its image is exactly the source
relative Lang system, because odd pullback identifies the constant
ambient cohomology as well. Removing the extra block determinants
from the source smooth, sharp open preserves its geometric image.
Hence the slice image \(G_0\subset G\) preserves this splitting
and projects surjectively onto the actual source image \(G_B\).

At the explicitly constructed composite common-wall point, every
outer character cycle remains bijective and the source has stable
loss one. The regular Lang hyperplane of the ambient full radical
is consequently \(H_0\oplus V_{\mathrm{nt}}\), with \(H_0\)
a proper hyperplane of \(V_0\). The strictly henselian comparison
of regular Lang lifts in
Proposition~\ref{prop:DE-Lang-deck-slice-quotient-eight} identifies
this hyperplane with the stated slice splitting; it does not
identify full generic fibres across the rank drop. Since \(K\)
lies in every full-radical fixed hyperplane, \(e(K)\subset H_0\).

Let \(W=e(K)\subset V_0\) at that base point. The projector
commutes with the slice image, and \(K\) is \(G\)-stable.
Thus \(W\) is stable under \(G_0\), and hence under its full
projected image \(G_B\). The inclusion in \(H_0\) makes \(W\)
proper. The source full-radical structure theorem says that its
canonical \(K_0\) is the unique maximal proper \(G_B\)-submodule.
Therefore \(W\subseteq K_0\), proving (DE-CR2). When \(K_0=0\)
this recovers the full-source conclusion, but no such vanishing
was used. The argument includes \(a=1\) and \(n=1\).
\end{proof}

\begin{theorem}[All intermediate residual spaces on one fibre]
\label{thm:DE-compatible-residual-module-sum}
\status{Candidate; complete proof supplied for independent audit.}
Let \(\pi:A\to B\) be an etale isogeny with
\(D=\ker\pi\simeq(\mathbf Z/\ell)^2\), where \(\ell\)
is an odd prime and \((B,\Theta_B)\) is nonordinary and
principally polarized. Put \(L=\pi^*\mathcal O_B(3\Theta_B)\).
For every line \(H\subset D\), let
\[
 A_H=A/H,\qquad \eta_H:A_H\to B,\qquad
 M_H=\eta_H^*\mathcal O_B(3\Theta_B),
 \tag{DE-CR3}
\]
and let \(K_H\) be the canonical residual module of the complete
source system \(|M_H|\). Write
\(m_H=9\ell-1-\dim K_H\).

There is a single common smooth, sharp cubic slice on which all
these spaces are identified with actual sub-local systems of
the target. At a geometric point of that slice,
\[
 K\subseteq\bigoplus_{|H|=\ell}K_H,
 \qquad
 \dim(V/K)\ge 8+\sum_{|H|=\ell}(m_H-8).
 \tag{DE-CR4}
\]
The displayed direct sum uses distinct actual deck sectors on
this same fibre, not unrelated identifications of generic spaces.
\end{theorem}
\begin{proof}
Use the common open \(S_D\) and connected intermediate opens
\(T_H\) constructed in
Proposition~\ref{prop:compatible-intermediate-common-slice}.
They are defined by source and target smoothness, sharpness and
the invertibility of the outer character cycles, and each \(T_H\)
contains the actual pullback of \(S_D\).
Over \(T_H\), the invariant part \(e_H\mathcal V\) is the
source relative local system. Its canonical source sub-local
system is denoted \(\mathcal K_H\); it is the restriction of
the source system defined on its full smooth, sharp open.

The preceding lemma gives on a nonempty open of \(T_H\) the
vanishing of the actual morphism
\[
 \mathcal K|_{T_H}\longrightarrow e_H\mathcal V|_{T_H}
           \longrightarrow e_H\mathcal V|_{T_H}/\mathcal K_H.
 \tag{DE-CR5}
\]
As a morphism of finite etale local systems on connected \(T_H\),
it vanishes everywhere if it vanishes on a nonempty open.
Consequently \(e_H\mathcal K\subseteq\mathcal K_H\) on
the same \(S_D\) for every line \(H\). This extension is the
zero-morphism argument of the compatible-intermediate theorem,
now with the actual source residual quotient as target.

For this further continuation, define \(R_H\) to be the open
in \(\mathbf P(H^0(B,3\Theta_B))\) where the degree-one curve
and its \(\eta_H\)-pullback are smooth and sharp and all
nontrivial \(\eta_H\)-character multiplication blocks are
invertible. The explicit central character construction makes
\(R_H\) nonempty and integral, and \(S_D\subset R_H\).
The pullback of the source \(\mathcal K_H\) and its central
averaging projector are therefore defined throughout \(R_H\).
The narrow deck-slice theorem for \(\eta_H\) says that their
composition is zero on a nonempty open of \(R_H\). Connectedness
and the same zero-morphism argument extend this to all of
\(R_H\), hence to \(S_D\). In particular no continuation is
made across a point where the intermediate pullback is not sharp.
Under the actual pullback to the target fibre, this central
projector is precisely \(e_D\) restricted to \(e_HV\).
Indeed the quotient deck group is \(D/H\), and averaging its
cosets after the odd \(H\)-average gives \(e_D\). Thus
\[
           K_H\subseteq e_HV\cap\ker e_D.
 \tag{DE-CR6}
\]
The same narrow slice theorem for \(\pi\) gives \(e_DK=0\)
on \(S_D\), with the identical connected-open justification.

Everything now takes place in the one actual
\(\mathbf F_2[D]\)-module \(V\). Define
\(p_H=e_H+e_D\). For distinct lines \(H,J\), the product
\(e_He_J=e_D\), since \(H+J=D\) and each element has a
unique such factorization. Together with (CCI3), this gives
\[
 p_H^2=p_H,\quad p_Hp_J=0\ (H\ne J),\quad
 e_Dp_H=0,\quad 1=e_D+\sum_{|H|=\ell}p_H.
 \tag{DE-CR7}
\]
Moreover \(p_HV=e_HV\cap\ker e_D\). The central summand
\(e_DV\) has dimension eight, while \(e_HV\) has dimension
\(9\ell-1\). Hence each \(p_HV\) has dimension
\(9(\ell-1)\), and these summands are pairwise disjoint.
This is a group-algebra calculation over \(\mathbf F_2\);
no character-splitting field is required.

For \(v\in K\), one has \(e_Dv=0\) and
\(p_Hv=e_Hv\in K_H\). Equation (DE-CR7) therefore gives
the first inclusion of (DE-CR4), with the direct sum justified
by (DE-CR6). Taking dimensions gives
\[
 \dim K\le\sum_H\dim K_H
           =\sum_H(9\ell-1-m_H).
\]
Since \(\dim V=9\ell^2-1\) and there are \(\ell+1\)
lines, rearranging this inequality gives exactly the second
part of (DE-CR4).
\end{proof}

\begin{corollary}[A quadratic quotient bound at a prime target degree]
\label{cor:DE-prime-compatible-wedge-bound}
\status{Candidate; depends on the separate complete wedge-bound candidate.}
In the theorem, the determinant-wedge bound gives
\[
 \dim(V/K)\ge3\ell^2-5\ell,
 \qquad \dim K\le6\ell^2+5\ell-1.
 \tag{DE-CR8}
\]
In particular, for every fixed nonordinary principally polarized
surface in degree \(\ell\), the main nonzero Lang component has
degree at least
\[
          2^{9\ell^2-1}-2^{6\ell^2+5\ell-1}.
 \tag{DE-CR9}
\]
There is also the sharper exact integer bound. Put
\(f=\operatorname{ord}_\ell(2)\). Then
\[
 \dim K\le (\ell+1)f
                 \left\lfloor\frac{6\ell-1}{f}\right\rfloor,
 \qquad
 \dim(V/K)\ge9\ell^2-1-(\ell+1)f
                 \left\lfloor\frac{6\ell-1}{f}\right\rfloor.
 \tag{DE-CR10}
\]
In the fixed-principal case this sharper exceptional bound can
likewise replace the exponent in the subtracted term of (DE-CR9).
At \(\ell=17\), one has \(f=8\), each source exceptional
dimension is at most 96, and the target quotient dimension is
at least 872 with exceptional dimension at most 1728, in total
rank 2600.
\end{corollary}
\begin{proof}
Each \(M_H\) has polarization type \((3,3\ell)\), so
Corollary~\ref{cor:DE-Lang-wedge-cyclic-principal-bounds} gives
\(m_H\ge3\ell\). Substitution into (DE-CR4) yields
\(8+(\ell+1)(3\ell-8)=3\ell^2-5\ell\), and its
complement in rank \(9\ell^2-1\) is the stated upper bound.
For a fixed principal target, take its exact outgoing
maximal-isotropic \(\ell\)-quotient and apply the already
proved main-orbit formula in that fixed-principal-system scope.
To prove the refinement, work at the same central geometric
point. Every deck translation of \(\eta_H\), whose group is
\(D/H\simeq C_\ell\), acts on the entire source universal
family and its smooth, sharp coefficient open. It fixes the
central section, because that section is pulled back from \(B\).
Functoriality of the fibre representation therefore makes its
linear action on the source Lang fibre normalize the source
geometric monodromy image. Since \(K_H\) is that image's unique
maximal proper submodule, it is invariant under these deck
translations. Thus it is an actual \(\mathbf F_2[C_\ell]\)-module.
Its averaging projector is zero by (DE-CR6).

The algebra \(\mathbf F_2[C_\ell]\) is semisimple because
\(\ell\) is odd. Every nontrivial simple module has dimension
\(f=\operatorname{ord}_\ell(2)\): the nonidentity \(\ell\)-th
roots have Frobenius orbits of exactly that length. Since the
trivial summand has been excluded, \(\dim K_H\) is a multiple
of \(f\). Its wedge bound \(\dim K_H\le6\ell-1\) therefore
improves to \(f\lfloor(6\ell-1)/f\rfloor\). Summing the
actual direct-sum inclusion (DE-CR4) proves (DE-CR10).
For \(\ell=17\), \(2^4\equiv-1\pmod {17}\) and
\(2^8\equiv1\pmod {17}\), while no proper divisor of eight
gives one. Hence \(f=8\), and the displayed numbers follow
by substitution. No finite experiment enters the proof.
\end{proof}

\begin{remark}[What the compatible residual bound leaves open]
\label{rem:DE-compatible-residual-module-scope}
\status{Candidate scope; no new unconditional full image.}
The first lemma compares actual residual modules and imposes no
source full-image hypothesis. The compatible theorem uses only
intermediate quotients of one fixed rank-two deck group and a
proved common smooth, sharp slice. It does not compare disjoint
eigenvector slices from unrelated theta groups.
The final corollary is conditional on completion of the separate
determinant-wedge proof audits. Even with that bound, the source
residual modules may be nonzero; the resulting target bound does
not prove \(K=0\) or full A.4. All proper submodule and local-system
claims use the previously proved full-radical and common-wall
arguments with their actual rank hypotheses. No new numerical
program or external source is needed or claimed.
\end{remark}
